% Options for packages loaded elsewhere
\PassOptionsToPackage{unicode}{hyperref}
\PassOptionsToPackage{hyphens}{url}
\PassOptionsToPackage{dvipsnames,svgnames,x11names}{xcolor}
%
\documentclass[
  11pt,
  a4paper,
]{ltjsarticle}
\usepackage{amsmath,amssymb}
\usepackage{iftex}
\ifPDFTeX
  \usepackage[T1]{fontenc}
  \usepackage[utf8]{inputenc}
  \usepackage{textcomp} % provide euro and other symbols
\else % if luatex or xetex
  \usepackage{unicode-math} % this also loads fontspec
  \defaultfontfeatures{Scale=MatchLowercase}
  \defaultfontfeatures[\rmfamily]{Ligatures=TeX,Scale=1}
\fi
\usepackage{lmodern}
\ifPDFTeX\else
  % xetex/luatex font selection
\fi
% Use upquote if available, for straight quotes in verbatim environments
\IfFileExists{upquote.sty}{\usepackage{upquote}}{}
\IfFileExists{microtype.sty}{% use microtype if available
  \usepackage[]{microtype}
  \UseMicrotypeSet[protrusion]{basicmath} % disable protrusion for tt fonts
}{}
\makeatletter
\@ifundefined{KOMAClassName}{% if non-KOMA class
  \IfFileExists{parskip.sty}{%
    \usepackage{parskip}
  }{% else
    \setlength{\parindent}{0pt}
    \setlength{\parskip}{6pt plus 2pt minus 1pt}}
}{% if KOMA class
  \KOMAoptions{parskip=half}}
\makeatother
\usepackage{xcolor}
\usepackage[margin=24mm]{geometry}
\setlength{\emergencystretch}{3em} % prevent overfull lines
\providecommand{\tightlist}{%
  \setlength{\itemsep}{0pt}\setlength{\parskip}{0pt}}
\setcounter{secnumdepth}{5}
\ifLuaTeX
  \usepackage{selnolig}  % disable illegal ligatures
\fi
\IfFileExists{bookmark.sty}{\usepackage{bookmark}}{\usepackage{hyperref}}
\IfFileExists{xurl.sty}{\usepackage{xurl}}{} % add URL line breaks if available
\urlstyle{same}
\hypersetup{
  colorlinks=true,
  linkcolor={blue},
  filecolor={Maroon},
  citecolor={Blue},
  urlcolor={blue},
  pdfcreator={LaTeX via pandoc}}

\author{}
\date{}

\begin{document}

\hypertarget{chapter-05-ux884cux5217ux5f0f}{%
\section{Chapter 05 --- 行列式}\label{chapter-05-ux884cux5217ux5f0f}}

行列式は公式展開が目立つため、線形代数で最も意味を見失いやすい概念の一つです。

\hypertarget{ux4e00ux3064ux306eux610fux5473ux306bux56faux5b9aux3059ux308b}{%
\subsection{一つの意味に固定する}\label{ux4e00ux3064ux306eux610fux5473ux306bux56faux5b9aux3059ux308b}}

まず

\[
\det A=\text{向き付き体積の倍率}
\]

と理解してください。

この一つの解釈から多くの性質が自然になります。

\begin{itemize}
\tightlist
\item
  \texttt{det\ A=0}: 体積が潰れた → 次元が落ちた → 非可逆
\item
  \texttt{det(AB)=det\ A\ det\ B}: 二つの変換の体積倍率は積
\item
  負の determinant: 向きが反転
\end{itemize}

\hypertarget{svd-ux3068ux306eux95a2ux4fc2}{%
\subsection{SVD との関係}\label{svd-ux3068ux306eux95a2ux4fc2}}

\[
|\det A|=\prod_i\sigma_i
\]

です。SVD で各直交軸を \texttt{σ\_i}
倍するので、体積はその積だけ変わります。

ここで「行列式が小さい」と「条件数が大きい」を混同しないことが重要です。全方向を一様に
\texttt{10\^{}\{-6\}} 倍する行列は determinant
は小さいですが、方向間の歪みがないため条件数は1です。

\hypertarget{ux5230ux9054ux57faux6e96}{%
\subsection{到達基準}\label{ux5230ux9054ux57faux6e96}}

余因子展開を知らなくても、行列式が0になることとランク落ち・非可逆性の関係を幾何学的に説明できる状態を目指します。

\end{document}
